\documentclass[aspectratio=169,10pt]{beamer}
\usetheme[numbering=fraction,progressbar=frametitle]{metropolis}
\usepackage{fontspec}
\setsansfont{Latin Modern Sans}
\usepackage{amsmath,amssymb,mathtools,graphicx,booktabs}
\usepackage{polyglossia}
\setdefaultlanguage{french}
\setbeamertemplate{navigation symbols}{}
\setbeamerfont{frametitle}{size=\large}
\setbeamerfont{normal text}{size=\normalsize}
\setlength{\parskip}{.45em}
\hypersetup{pdfauthor={},pdfsubject={Lecture slides},pdfkeywords={},pdftitle={Limites de suites}}
\providecommand{\N}{\mathbb{N}}
\providecommand{\Z}{\mathbb{Z}}
\providecommand{\D}{\mathbb{D}}
\providecommand{\Q}{\mathbb{Q}}
\providecommand{\R}{\mathbb{R}}
\providecommand{\C}{\mathbb{C}}
\providecommand{\vect}[1]{\overrightarrow{#1}}
\providecommand{\vi}{\vec{\imath}}
\providecommand{\vj}{\vec{\jmath}}
\providecommand{\vk}{\vec{k}}
\providecommand{\norm}[1]{\left\lVert #1\right\rVert}
\providecommand{\vc}[2]{\begin{pmatrix}#1\\#2\end{pmatrix}}
\providecommand{\deter}[2]{\det\left(#1,#2\right)}
\providecommand{\abs}[1]{\left\lvert #1\right\rvert}
\providecommand{\dd}{\,\mathrm{d}}
\providecommand{\e}{\mathrm{e}}
\providecommand{\Card}{\operatorname{Card}}
\providecommand{\E}{\operatorname{E}}
\providecommand{\V}{\operatorname{V}}
\providecommand{\degre}{^{\circ}}
\providecommand{\un}[1]{\,\mathrm{#1}}
\providecommand{\sci}[2]{#1\times10^{#2}}
\providecommand{\pgcd}{\operatorname{PGCD}}
\providecommand{\ppcm}{\operatorname{PPCM}}
\title{Limites de suites}
\subtitle{Cours et exemples}
\date{}
\author{}
\begin{document}
\begin{frame}[plain]
\vfill
{\usebeamerfont{title}\inserttitle\par}
\medskip
{\usebeamerfont{subtitle}\insertsubtitle\par}
\medskip
{\usebeamercolor[fg]{progress bar}\rule{\textwidth}{.5pt}}
\vfill
\end{frame}
\begin{frame}{Objectifs}
Distinguer convergence et divergence ; calculer une limite ; utiliser comparaison, gendarmes et convergence monotone.
\end{frame}
\begin{frame}{Notes de cours}
\(u_n\to\ell\in\R\) signifie que tout intervalle ouvert contenant \(\ell\) contient tous les termes à partir d'un certain rang. On dit que \((u_n)\) converge vers \(\ell\).
\(u_n\to+\infty\) signifie que pour tout réel \(A\), il existe un rang \(N\) tel que \(n\ge N\Rightarrow u_n\ge A\). \(u_n\to-\infty\) signifie que pour tout réel \(A\), il existe \(N\) tel que \(n\ge N\Rightarrow u_n\le A\).
\end{frame}
\begin{frame}{Notes de cours}
Une suite peut n'avoir aucune limite, même infinie : \((-1)^n\) oscille entre \(-1\) et \(1\). Toute suite convergente est bornée, mais la réciproque est fausse.
\end{frame}
\begin{frame}{Opérations et comparaisons}
Pour deux suites convergentes de limites \(a,b\), la somme tend vers \(a+b\), le produit vers \(ab\) et le quotient vers \(\dfrac{a}{b}\) si \(b\ne0\). Les règles avec l'infini s'utilisent en contrôlant les signes. Les formes \(+\infty-\infty\), \(0\times\infty\), \(\dfrac{\infty}{\infty}\), \(\dfrac{0}{0}\) ne permettent pas de conclure directement.
\end{frame}
\begin{frame}{Opérations et comparaisons}
\(\dfrac{3n^2-n+1}{2n^2+4}=\dfrac{3-\dfrac{1}{n}+\dfrac{1}{n^2}}{2+\dfrac{4}{n^2}}\to\dfrac{3}{2}\). Pour \(\sqrt{n^2+n}-n\), multiplier par la quantité conjuguée donne \(\dfrac{1}{\sqrt{1+\dfrac{1}{n}}+1}\to\dfrac{1}{2}\).
\end{frame}
\begin{frame}{Opérations et comparaisons}
Si \(u_n\le v_n\) à partir d'un certain rang et \(u_n\to+\infty\), alors \(v_n\to+\infty\). Si \(a_n\le u_n\le b_n\) à partir d'un certain rang et \(a_n,b_n\to\ell\), alors \(u_n\to\ell\) (gendarmes).
Toute suite croissante et majorée converge. Toute suite décroissante et minorée converge. Ce théorème est admis.
\end{frame}
\begin{frame}{Opérations et comparaisons}
Pour tout \(A\), les termes de \(u\) dépassent \(A\) au-delà d'un rang \(N_1\). L'inégalité \(v_n\ge u_n\) vaut au-delà d'un rang \(N_2\). Pour \(n\ge\max(N_1,N_2)\), \(v_n\ge A\).
\end{frame}
\begin{frame}{Opérations et comparaisons}
Pour tout \(A\), la non-majoration fournit \(N\) tel que \(u_N>A\). Par croissance, \(n\ge N\Rightarrow u_n\ge u_N>A\), d'où \(u_n\to+\infty\). La non-majoration seule ne suffit pas.
\end{frame}
\begin{frame}{Suites géométriques}
\[q^n\longrightarrow\begin{cases}0&\text{si }-1<q<1,\\1&\text{si }q=1,\\+\infty&\text{si }q>1.\end{cases}\]
Si \(q\le-1\), \((q^n)\) n'a pas de limite : alternance pour \(q=-1\) ; sous-suites paire et impaire de signes opposés et de modules non bornés pour \(q<-1\).
\end{frame}
\begin{frame}{Suites géométriques}
Fixons \(a\ge0\). Démontrons \(P(n)\) : « \((1+a)^n\ge1+na\) ». Au rang \(0\), les deux membres valent \(1\).
Fixons \(n\ge0\) et supposons \(P(n)\) vraie. Comme \(1+a>0\),
\[(1+a)^{n+1}\ge(1+na)(1+a)=1+(n+1)a+na^2\ge1+(n+1)a.\]
C'est \(P(n+1)\) car \(na^2\ge0\). Par récurrence, l'inégalité est vraie pour tout \(n\ge0\). Si \(q>1\), prendre \(a=q-1>0\) et comparer à \(1+na\). Si \(0<q<1\), appliquer ce résultat à \(\dfrac{1}{q}\), puis prendre les inverses. Pour \(-1<q<0\), \(|q^n|=|q|^n\to0\) et les gendarmes concluent. Le cas \(q=0\) se traite à partir du rang \(1\).
\end{frame}
\begin{frame}{Suites géométriques}
Pour \(u_{n+1}=f(u_n)\), établir une borne et la monotonie pour prouver la convergence. Si \(f\) est continue en la limite \(\ell\), le passage à la limite donne \(\ell=f(\ell)\) (continuité étudiée au chapitre 4). Résoudre cette équation ne prouve jamais, à lui seul, que la suite converge.
\end{frame}
\begin{frame}{Suites géométriques}
\begin{columns}[T,onlytextwidth]
\begin{column}{.57\textwidth}
Les points se rapprochent de \(2\) sans l'atteindre : \(u_n=2+\dfrac{3}{n+1}\).
\end{column}
\begin{column}{.4\textwidth}
\centering
\includegraphics[width=\linewidth,height=4.5cm,keepaspectratio]{figures/terminale-03-1.pdf}
\end{column}
\end{columns}
\end{frame}
\begin{frame}{Activité de recherche}
\textbf{Une limite n'est pas un dernier terme - 30 min.}

On considère la suite \((u_n)\) définie sur \(\N\) par \(u_n=2+\dfrac3{n+1}\).
\end{frame}
\begin{frame}{Activité de recherche (2)}
\begin{enumerate}
\setcounter{enumi}{0}\item Calculer \(u_0\), \(u_9\) et \(u_{99}\). Existe-t-il un entier \(n\) tel que \(u_n=2\) ? Justifier.

\item Déterminer le plus petit entier \(N\) tel que, pour tout entier \(n\ge N\), \(|u_n-2|<0{,}01\).

\item Soit \(\varepsilon>0\). Donner, en fonction de \(\varepsilon\), un entier \(N\) tel que \(|u_n-2|<\varepsilon\) pour tout \(n\ge N\).

\item On définit sur \(\N\) la suite \((v_n)\) par \(v_n=2+\dfrac{(-1)^n}{n+1}\). Démontrer qu'elle converge vers \(2\) et qu'elle n'est pas monotone. Que montre cet exemple ?

\end{enumerate}
\end{frame}
\end{document}
